\documentclass[aspectratio=169, 12pt]{beamer}

% ── Theme ──────────────────────────────────────────────
\usetheme{default}
\usecolortheme{default}
\setbeamertemplate{navigation symbols}{}
\setbeamertemplate{footline}{}

% ── Packages ───────────────────────────────────────────
\usepackage[utf8]{inputenc}
\usepackage[T1]{fontenc}
\usepackage[spanish]{babel}
\usepackage{lmodern}
\usepackage{tikz}
\usepackage{xcolor}
\usepackage{fontawesome5}
\usepackage{tcolorbox}
\usepackage{booktabs}
\usepackage{graphicx}
\usepackage{array}
\usepackage{calc}
\tcbuselibrary{skins}

% ── Colors ─────────────────────────────────────────────
\definecolor{navy}{RGB}{13,53,89}        % #0D3559
\definecolor{navydark}{RGB}{26,42,74}    % #1A2A4A
\definecolor{azul}{RGB}{58,171,227}      % #3AABE3
\definecolor{verde}{RGB}{75,173,78}      % #4BAD4E
\definecolor{rojo}{RGB}{232,48,42}       % #E8302A
\definecolor{naranja}{RGB}{244,121,32}   % #F47920
\definecolor{amarillo}{RGB}{249,194,0}   % #F9C200
\definecolor{morado}{RGB}{139,78,166}    % #8B4EA6
\definecolor{bglight}{RGB}{244,246,251}  % #F4F6FB
\definecolor{textgray}{RGB}{102,102,119} % #666677

% ── Beamer colors ──────────────────────────────────────
\setbeamercolor{background canvas}{bg=bglight}
\setbeamercolor{frametitle}{fg=white, bg=navy}
\setbeamercolor{title}{fg=white}
\setbeamercolor{subtitle}{fg=azul}
\setbeamercolor{normal text}{fg=navydark}
\setbeamercolor{itemize item}{fg=azul}
\setbeamercolor{itemize subitem}{fg=textgray}
\setbeamercolor{block title}{fg=white, bg=navy}
\setbeamercolor{block body}{bg=white}

% ── Fonts ──────────────────────────────────────────────
\setbeamerfont{title}{size=\LARGE, series=\bfseries}
\setbeamerfont{subtitle}{size=\large}
\setbeamerfont{frametitle}{size=\large, series=\bfseries}
\setbeamerfont{itemize/enumerate body}{size=\normalsize}

% ── Frame title with rainbow bar ───────────────────────
\setbeamertemplate{frametitle}{%
  \vskip-1pt
  \begin{beamercolorbox}[wd=\paperwidth, ht=0.5cm, dp=0cm]{frametitle}%
    \begin{tikzpicture}
      \fill[rojo]    (0,0)                     rectangle (0.1667\paperwidth, 0.08cm);
      \fill[naranja] (0.1667\paperwidth,0)     rectangle (0.3333\paperwidth, 0.08cm);
      \fill[amarillo](0.3333\paperwidth,0)     rectangle (0.5\paperwidth,    0.08cm);
      \fill[verde]   (0.5\paperwidth,0)        rectangle (0.6667\paperwidth, 0.08cm);
      \fill[azul]    (0.6667\paperwidth,0)     rectangle (0.8333\paperwidth, 0.08cm);
      \fill[morado]  (0.8333\paperwidth,0)     rectangle (\paperwidth,       0.08cm);
    \end{tikzpicture}
  \end{beamercolorbox}%
  \begin{beamercolorbox}[wd=\paperwidth, ht=0.7cm, dp=0.2cm, leftskip=0.4cm]{frametitle}%
    \usebeamerfont{frametitle}\insertframetitle
  \end{beamercolorbox}%
}

% ── Itemize bullets ────────────────────────────────────
\setbeamertemplate{itemize item}{\textcolor{azul}{\textbullet}}
\setbeamertemplate{itemize subitem}{\textcolor{textgray}{--}}

% ── Helper: colored pill ───────────────────────────────
\newcommand{\pill}[2]{%
  \tikz[baseline=-0.6ex]\node[fill=#1, text=white, rounded corners=8pt,
    inner xsep=6pt, inner ysep=2pt, font=\footnotesize\bfseries]{#2};%
}
\newcommand{\pilloutline}[2]{%
  \tikz[baseline=-0.6ex]\node[draw=#1, text=#1, rounded corners=8pt,
    inner xsep=6pt, inner ysep=2pt, font=\footnotesize\bfseries, line width=1.2pt]{#2};%
}

% ── Helper: stat box ───────────────────────────────────
\newcommand{\statbox}[3]{%
  \begin{tcolorbox}[colback=white, colframe=#1, arc=8pt,
    boxrule=2pt, width=0.28\textwidth, halign=center, valign=center,
    top=6pt, bottom=6pt]
    {\color{#1}\fontsize{22}{26}\selectfont\bfseries #2}\\[2pt]
    {\color{textgray}\footnotesize #3}
  \end{tcolorbox}%
}

% ══════════════════════════════════════════════════════
%   TITLE PAGE
% ══════════════════════════════════════════════════════

\title{\textbf{AulaPresente}}
\subtitle{Sistematización de asistencia escolar en Bogotá}
\author{Oscar M. Díaz-Botia\\
  {\small Paris School of Economics}\\[4pt]
  {\small \textcolor{azul}{oscarmdiazb@gmail.com}}}
\date{\small Bogotá, 2025}

\begin{document}

% ── Slide 1: Title ─────────────────────────────────────
{
\setbeamertemplate{frametitle}{}
\setbeamercolor{background canvas}{bg=navy}
\begin{frame}
  \vspace{0.3cm}
  % Rainbow bar
  \begin{tikzpicture}[remember picture, overlay]
    \fill[rojo]    (0,0)                     rectangle (0.1667\paperwidth, 0.12cm);
    \fill[naranja] (0.1667\paperwidth,0)     rectangle (0.3333\paperwidth, 0.12cm);
    \fill[amarillo](0.3333\paperwidth,0)     rectangle (0.5\paperwidth,    0.12cm);
    \fill[verde]   (0.5\paperwidth,0)        rectangle (0.6667\paperwidth, 0.12cm);
    \fill[azul]    (0.6667\paperwidth,0)     rectangle (0.8333\paperwidth, 0.12cm);
    \fill[morado]  (0.8333\paperwidth,0)     rectangle (\paperwidth,       0.12cm);
  \end{tikzpicture}
  \vspace{0.6cm}
  \begin{center}
    {\color{white}\fontsize{36}{42}\selectfont\bfseries AulaPresente}\\[8pt]
    {\color{azul}\large Sistematización de asistencia escolar en Bogotá}\\[20pt]
    \textcolor{white!60}{\rule{6cm}{0.5pt}}\\[14pt]
    {\color{white!80}\small Oscar M. Díaz-Botia \quad · \quad Paris School of Economics}\\[4pt]
    {\color{white!60}\footnotesize Bogotá, 2025}
  \end{center}
\end{frame}
}

% ── Slide 2: El problema ───────────────────────────────
\begin{frame}{El problema: un punto ciego en la política educativa}
\vspace{0.2cm}
\begin{columns}[T]
  \begin{column}{0.55\textwidth}
    \textbf{\textcolor{navy}{Bogotá no tiene datos de asistencia sistematizados.}}\\[8pt]
    Los colegios registran la asistencia de formas dispares:
    papel, Excel, cuadernos propios. Los datos \textbf{no se consolidan
    ni se analizan} a nivel de ciudad.\\[12pt]
    Esto impide:
    \begin{itemize}\setlength\itemsep{4pt}
      \item Medir la magnitud real del ausentismo
      \item Identificar patrones por localidad o grado
      \item Detectar estudiantes en riesgo de deserción
      \item Tomar decisiones de política basadas en evidencia
    \end{itemize}
  \end{column}
  \begin{column}{0.42\textwidth}
    \begin{tcolorbox}[colback=rojo!8, colframe=rojo, arc=10pt,
      boxrule=1.5pt, title={\color{white}\bfseries Magnitud del problema},
      coltitle=white, colbacktitle=rojo]
      \centering
      {\color{rojo}\fontsize{28}{32}\selectfont\bfseries ?}\\[4pt]
      {\small\color{textgray} No existen cifras consolidadas de ausentismo escolar en Bogotá}\\[8pt]
      {\color{rojo}\fontsize{20}{24}\selectfont\bfseries 700k+}\\[4pt]
      {\small\color{textgray} estudiantes en colegios públicos de Bogotá}
    \end{tcolorbox}
  \end{column}
\end{columns}
\end{frame}

% ── Slide 3: La solución ───────────────────────────────
\begin{frame}{La solución: AulaPresente}
\vspace{0.1cm}
\begin{center}
  {\large Una app web simple para registrar asistencia desde el celular,}\\[2pt]
  {\large diseñada para docentes y monitores en contextos reales.}
\end{center}
\vspace{0.3cm}
\begin{columns}[T]
  \begin{column}{0.3\textwidth}
    \begin{tcolorbox}[colback=azul!8, colframe=azul, arc=10pt,
      boxrule=1.5pt, halign=center, top=8pt, bottom=8pt]
      {\color{azul}\Large\bfseries 1}\\[4pt]
      {\bfseries Elige el grupo}\\[4pt]
      {\small\color{textgray} Selecciona la clase, la fecha y el tipo de sesión}
    \end{tcolorbox}
  \end{column}
  \begin{column}{0.3\textwidth}
    \begin{tcolorbox}[colback=naranja!8, colframe=naranja, arc=10pt,
      boxrule=1.5pt, halign=center, top=8pt, bottom=8pt]
      {\color{naranja}\Large\bfseries 2}\\[4pt]
      {\bfseries Marca ausentes}\\[4pt]
      {\small\color{textgray} Toca el nombre del estudiante para marcarlo ausente}
    \end{tcolorbox}
  \end{column}
  \begin{column}{0.3\textwidth}
    \begin{tcolorbox}[colback=verde!8, colframe=verde, arc=10pt,
      boxrule=1.5pt, halign=center, top=8pt, bottom=8pt]
      {\color{verde}\Large\bfseries 3}\\[4pt]
      {\bfseries Guarda y exporta}\\[4pt]
      {\small\color{textgray} Registro guardado. Exportable como Excel en un toque}
    \end{tcolorbox}
  \end{column}
\end{columns}
\vspace{0.3cm}
\begin{center}
  \pill{navy}{Sin instalación} \quad
  \pill{azul}{Funciona sin internet} \quad
  \pill{verde}{Cualquier celular} \quad
  \pill{morado}{Compatible con SIMAT}
\end{center}
\end{frame}

% ── Slide 4: Características clave ────────────────────
\begin{frame}{Características clave}
\vspace{0.15cm}
\begin{columns}[T]
  \begin{column}{0.48\textwidth}
    \textcolor{navy}{\bfseries Para el usuario (docente / monitor)}
    \begin{itemize}\setlength\itemsep{5pt}
      \item Lista de estudiantes con N.\textdegree{} SIMAT, ordenada alfabéticamente
      \item Registro por fecha, tipo de sesión y franja horaria
      \item Historial completo de asistencia por clase
      \item Exportación a CSV / Excel con un toque
      \item Perfil del registrador adjunto a cada registro
    \end{itemize}
  \end{column}
  \begin{column}{0.48\textwidth}
    \textcolor{navy}{\bfseries Para la institución / SED}
    \begin{itemize}\setlength\itemsep{5pt}
      \item Datos estructurados y estandarizados
      \item Identificación por N.\textdegree{} SIMAT (cruce con SIMAT oficial)
      \item Registro del facilitador: nombre, rol, institución
      \item Escalable: puede conectarse a servidor SED en fase 2
      \item Sin costo de licencias (tecnología abierta)
    \end{itemize}
  \end{column}
\end{columns}
\vspace{0.25cm}
\begin{tcolorbox}[colback=verde!6, colframe=verde, arc=8pt, boxrule=1pt,
  left=6pt, right=6pt, top=4pt, bottom=4pt]
  \textcolor{verde}{\bfseries Privacidad:} Los datos se almacenan en el dispositivo del usuario. Nadie más tiene acceso a menos que el usuario exporte el archivo.
\end{tcolorbox}
\end{frame}

% ── Slide 5: El piloto ─────────────────────────────────
\begin{frame}{El piloto: monitores de movilidad escolar}
\vspace{0.2cm}
\begin{columns}[T]
  \begin{column}{0.55\textwidth}
    \textbf{\textcolor{navy}{¿Por qué los monitores de movilidad escolar?}}\\[8pt]
    \begin{itemize}\setlength\itemsep{6pt}
      \item Ya tienen presencia sistemática en los colegios
      \item Sus rutinas de trabajo incluyen registro de estudiantes
      \item Son usuarios naturales de una app de este tipo
      \item Permiten validar la herramienta en condiciones reales
      \item Su feedback informará las mejoras antes de escalar
    \end{itemize}
    \vspace{10pt}
    \textbf{\textcolor{navy}{Producto del piloto}}\\[6pt]
    \begin{itemize}\setlength\itemsep{4pt}
      \item Datos de asistencia estructurados y exportables
      \item Reporte de usabilidad y adopción
      \item Hoja de ruta para la fase de integración con SED
    \end{itemize}
  \end{column}
  \begin{column}{0.42\textwidth}
    \begin{tcolorbox}[colback=azul!6, colframe=azul, arc=10pt,
      boxrule=1.5pt, title={\color{white}\bfseries Acceso al piloto},
      coltitle=white, colbacktitle=azul, halign=center]
      \vspace{4pt}
      {\small Disponible hoy en:}\\[4pt]
      {\footnotesize\ttfamily\color{navy}
        oscarmdiazb.github.io/\\asistencia-escolar/}\\[8pt]
      {\small\color{textgray} Se abre en Chrome o Safari.\\
       No requiere instalación.\\
       Funciona en cualquier celular.}
      \vspace{4pt}
    \end{tcolorbox}
    \vspace{8pt}
    \begin{tcolorbox}[colback=amarillo!15, colframe=amarillo, arc=8pt,
      boxrule=1pt, halign=center]
      {\small\bfseries También se usará en el marco del proyecto}\\[2pt]
      {\small\color{textgray} \textit{Aulas en Re-Conexión} (RCT, Paris School of Economics)}
    \end{tcolorbox}
  \end{column}
\end{columns}
\end{frame}

% ── Slide 6: Hoja de ruta ──────────────────────────────
\begin{frame}{Hoja de ruta hacia la sistematización ciudad}
\vspace{0.3cm}
\begin{tikzpicture}[node distance=0cm]

  % Phase boxes
  \node[fill=verde, text=white, rounded corners=8pt,
        minimum width=3.8cm, minimum height=1.8cm,
        text width=3.4cm, align=center, font=\small\bfseries] (f1) at (0,0)
        {FASE 1\\[4pt]\normalfont Piloto\\monitores SED\\(2025)};

  \node[fill=azul, text=white, rounded corners=8pt,
        minimum width=3.8cm, minimum height=1.8cm,
        text width=3.4cm, align=center, font=\small\bfseries] (f2) at (4.4,0)
        {FASE 2\\[4pt]\normalfont Integración\\servidor SED\\(2026)};

  \node[fill=navy, text=white, rounded corners=8pt,
        minimum width=3.8cm, minimum height=1.8cm,
        text width=3.4cm, align=center, font=\small\bfseries] (f3) at (8.8,0)
        {FASE 3\\[4pt]\normalfont Escala ciudad\\todos los colegios\\(2027+)};

  % Arrows
  \draw[->, line width=2pt, color=textgray] (f1.east) -- (f2.west);
  \draw[->, line width=2pt, color=textgray] (f2.east) -- (f3.west);

  % Labels below
  \node[text width=3.4cm, align=center, font=\footnotesize\color{textgray}]
    at (0,-1.4) {App web · localStorage · GitHub Pages};
  \node[text width=3.4cm, align=center, font=\footnotesize\color{textgray}]
    at (4.4,-1.4) {Nube / API SED · datos en tiempo real};
  \node[text width=3.4cm, align=center, font=\footnotesize\color{textgray}]
    at (8.8,-1.4) {Integración SIMAT · app nativa · analítica};

\end{tikzpicture}
\vspace{0.4cm}
\begin{tcolorbox}[colback=navy!6, colframe=navy, arc=8pt, boxrule=1pt,
  left=8pt, right=8pt, top=5pt, bottom=5pt]
  \textbf{Resultado final:} la SED cuenta con datos de asistencia estandarizados para los
  700.000+ estudiantes de colegios públicos de Bogotá, en tiempo real, cruzables con SIMAT.
\end{tcolorbox}
\end{frame}

% ── Slide 7: Lo que pedimos ────────────────────────────
\begin{frame}{Lo que necesitamos para avanzar}
\vspace{0.3cm}
\begin{columns}[T]
  \begin{column}{0.48\textwidth}
    \begin{tcolorbox}[colback=white, colframe=verde, arc=10pt,
      boxrule=2pt, title={\color{white}\bfseries Para el piloto (Fase 1)},
      coltitle=white, colbacktitle=verde]
      \begin{itemize}\setlength\itemsep{5pt}
        \item Acceso a los monitores de movilidad escolar para el piloto
        \item Aval institucional de la SED para el uso experimental
        \item Retroalimentación de los usuarios durante el piloto
      \end{itemize}
    \end{tcolorbox}
  \end{column}
  \begin{column}{0.48\textwidth}
    \begin{tcolorbox}[colback=white, colframe=azul, arc=10pt,
      boxrule=2pt, title={\color{white}\bfseries Para la integración (Fase 2)},
      coltitle=white, colbacktitle=azul]
      \begin{itemize}\setlength\itemsep{5pt}
        \item Coordinación con la Oficina de TI de la SED
        \item Definición de estándares de datos y acceso a SIMAT
        \item Decisión sobre arquitectura de almacenamiento (nube SED vs. Firebase)
      \end{itemize}
    \end{tcolorbox}
  \end{column}
\end{columns}
\vspace{0.35cm}
\begin{center}
  \textcolor{textgray}{\small El desarrollo del app no tiene costo.
  Está listo para usarse hoy.\\
  La inversión requerida en Fase 1 es \textbf{solo tiempo y coordinación}.}
\end{center}
\end{frame}

% ── Slide 8: Cierre ─────────────────────────────────────
{
\setbeamertemplate{frametitle}{}
\setbeamercolor{background canvas}{bg=navy}
\begin{frame}
  \begin{tikzpicture}[remember picture, overlay]
    \fill[rojo]    (0,0)                     rectangle (0.1667\paperwidth, 0.12cm);
    \fill[naranja] (0.1667\paperwidth,0)     rectangle (0.3333\paperwidth, 0.12cm);
    \fill[amarillo](0.3333\paperwidth,0)     rectangle (0.5\paperwidth,    0.12cm);
    \fill[verde]   (0.5\paperwidth,0)        rectangle (0.6667\paperwidth, 0.12cm);
    \fill[azul]    (0.6667\paperwidth,0)     rectangle (0.8333\paperwidth, 0.12cm);
    \fill[morado]  (0.8333\paperwidth,0)     rectangle (\paperwidth,       0.12cm);
  \end{tikzpicture}
  \vspace{0.8cm}
  \begin{center}
    {\color{white}\Large\bfseries AulaPresente}\\[10pt]
    {\color{white!80}\normalsize
      Una herramienta simple con un impacto potencial enorme:\\[4pt]
      por primera vez, Bogotá podría saber quiénes van al colegio.}\\[24pt]
    \textcolor{white!40}{\rule{5cm}{0.4pt}}\\[14pt]
    {\color{azul}\small Prueba el app hoy:}\\[4pt]
    {\color{white}\footnotesize\ttfamily oscarmdiazb.github.io/asistencia-escolar/}\\[16pt]
    {\color{white!70}\small Oscar M. Díaz-Botia \quad · \quad oscarmdiazb@gmail.com}
  \end{center}
\end{frame}
}

\end{document}
